\documentclass[prd,aps,showpacs,nofootinbib,superscriptaddress,amsmath,amssymb]{revtex4}
\usepackage{amssymb,graphicx}
\usepackage{mathrsfs}
\usepackage[usenames,dvipsnames]{color}
\usepackage[bookmarks]{hyperref}
\usepackage{multirow, array}
\usepackage{float}

\newcommand{\Natural}{\mathbb{N}}
\newcommand{\Integer}{\mathbb{Z}}
\newcommand{\Real}{\mathbb{R}}
\newcommand{\Complex}{\mathbb{C}}
\newcommand{\re}{\mbox{Re}}
\newcommand{\im}{\mbox{Im}}
\newcommand{\supp}{\mbox{supp}\,}

% Comments of the authors. Remove before submission.
\newcommand{\EJ}[1]{{\color{red}~\textsf{[(Erik) #1]}}}
\newcommand{\OS}[1]{{\color{blue}~\textsf{[(Olivier) #1]}}}


%%%%%%%%%%%%%%%%%%%%%%%%%%%%%%%
\begin{document}
%%%%%%%%%%%%%%%%%%%%%%%%%%%%%%%


\title{Phase space mixing in the presence of self-gravity: spherically symmetric Vlasov-Poisson evolutions with fixed and dispersed angular momentum}

\author{Erik Jim\'enez-V\'azquez}
\affiliation{\EJ{Author list and affiliations to be completed.}}


\begin{abstract}
\EJ{Provisional; to be rewritten when sections~\ref{Sec:FixedL} and~\ref{Sec:Dispersion} are final.}
This article is devoted to the numerical study of a spherically symmetric, collisionless kinetic gas which is trapped in an external central potential and whose self-gravity is taken into account. In the absence of self-gravity, such a gas relaxes in time to a stationary state through phase space mixing. We analyze how this process is modified by the self-gravity of the gas, considering both the reduced model in which all gas particles have the same magnitude of the angular momentum and the case in which the angular momenta are distributed over an interval. Our study is based on a particle-in-cell code, on a numerical solution of the linearized equations in action-angle variables, and on a spectral criterion which characterizes the existence of oscillating modes. We find that small perturbations of a self-consistent equilibrium are damped as long as the mass of the gas lies below a threshold, whereas above this threshold they give rise to an oscillation which is not damped, and we show that the threshold and the frequency of the oscillation are accurately predicted by the spectral criterion. Finally, we show that a dispersion in the angular momentum suppresses the oscillations which are present in the reduced model.
\end{abstract}

\date{\today}

\pacs{05.20.Dd, 98.10.+z, 02.60.Cb}
% 05.20.Dd: Kinetic theory
% 98.10.+z: Stellar dynamics and kinematics
% 02.60.Cb: Numerical simulation; solution of equations

\maketitle


%%%%%%%%%%%%%%%%%%%%%%%%%%%%%%%%%%%%%%%%%%%%%
\section{Introduction}
\label{Sec:Introduction}
%%%%%%%%%%%%%%%%%%%%%%%%%%%%%%%%%%%%%%%%%%%%%

Collisionless kinetic gases provide the standard description for systems consisting of a large number of particles which interact with each other only through the mean field they generate. Prominent examples are the stars in a galaxy or in a globular cluster, and the dark matter particles in a galactic halo~\cite{BinneyTremaine-Book}. The state of such a gas at time $t$ is described by its one-particle distribution function (DF) $f(t,x,p)$, a nonnegative function on the one-particle phase space which determines the number of gas particles in a small phase space volume about the point $(x,p)$, and whose evolution is governed by the Vlasov equation.\footnote{The Vlasov equation is also often called the collisionless Boltzmann equation in the literature.} In the Newtonian gravitational case, the Vlasov equation is coupled to Poisson's equation for the gravitational potential generated by the gas, which yields the Vlasov-Poisson system, see Ref.~\cite{gR07} for a review of its mathematical properties. Its general relativistic counterpart is the Vlasov-Einstein system~\cite{hA11}.

It is instructive to compare the dynamics of a self-gravitating collisionless gas with the one of a pressureless fluid (dust), which is obtained in the limit in which the velocity dispersion of the gas vanishes. The gravitational collapse of a spherical dust cloud leads to the formation of a singularity in finite time at which the density diverges, both in the Newtonian theory and in general relativity~\cite{jOhS39}; in the latter case, part of this singularity may even be visible to distant observers~\cite{dC84,nOoS11}. In contrast to this, the velocity dispersion of a kinetic gas prevents the particles from focusing. In fact, the Vlasov-Poisson system admits global classical solutions for general initial data~\cite{kP92,pLbP91,jS91}, while for the spherically symmetric Vlasov-Einstein system global solutions exist for small initial data~\cite{gRaR92} and sufficiently concentrated data collapse to a black hole~\cite{hAmKgR11}, the threshold between these two outcomes having been explored numerically in Refs.~\cite{gRaRjS98,Olabarrieta:2001wy,hAgR06,aAmC14}. Therefore, for a Newtonian collisionless gas the relevant question is not whether or not a singularity forms, but rather how the gas behaves for large times and, in particular, whether or not it relaxes to a stationary state.

A fundamental mechanism leading to such a relaxation is phase space mixing~\cite{dL62,dL67}. To explain it, consider first a gas whose self-gravity is neglected and which is trapped in the potential well of an external potential. If the one-particle Hamiltonian is integrable, each gas particle is confined to an invariant torus, and in terms of action-angle variables $(Q,J)$ the solution of the Vlasov equation is $f(t,Q,J) = F(Q - \Omega(J)t,J)$, where $F$ describes the initial DF and $\Omega(J)$ denotes the frequencies associated with the torus labeled by $J$. Since these frequencies vary from one torus to the other, the DF is stretched along the angle directions and develops structures on finer and finer scales as time increases. As a consequence, $f$ does not converge pointwise as $t\to\infty$. However, the macroscopic observables which are obtained by integrating $f$ against a test function do converge, and their limit is the one associated with the average of the initial DF over the angle variables, which depends only on integrals of motion. A mathematically precise formulation of this statement and its proof for rotationally symmetric potentials in $d = 1,2,3$ space dimensions satisfying a suitable non-degeneracy condition were given by Rioseco and Sarbach~\cite{pRoS20} (see also Refs.~\cite{sTmHdL86,cM19}). The results in Ref.~\cite{pRoS20} include the case of a spherically symmetric gas configuration in which all the gas particles have the same value of the total angular momentum, which behaves like an effective one-dimensional system, and they explain the relaxation to a virialized state which had been observed in the numerical simulations of Ref.~\cite{pDeJmAeMdN17}. For the analogous phenomenon in the exterior of a black hole we refer the reader to Refs.~\cite{pRoS18,pRoS24}, and for estimates on the rate at which the observables converge to Refs.~\cite{mMpRhV22,sCjL22,mHgRmScS24}.

The generalization to the self-gravitating case, which was left open in Ref.~\cite{pRoS20}, is considerably more involved, since the potential now depends on the DF, such that a perturbation of the latter induces a perturbation of the potential which, in turn, acts back on the gas. In plasma physics, where the particles interact through their mean electric field, the corresponding collective effect has been known since the work of Landau~\cite{lL46}: for a gas of charged particles close to a spatially homogeneous equilibrium the electric field decays in time while the DF undergoes mixing. This phenomenon, known as Landau damping, has been established rigorously at the nonlinear level by Mouhot and Villani~\cite{cMcV11}. The gravitational case differs from the electromagnetic one in two important aspects. First, the interaction is attractive, such that the collective response of the gas tends to reinforce a perturbation instead of screening it. Second, a self-gravitating gas configuration of finite mass does not possess spatially homogeneous equilibria. Instead, its steady states are spatially inhomogeneous, the gas particles are trapped in the potential well, and the frequencies of their radial motion lie in an interval which, in the situations considered in these notes, is bounded and bounded away from zero. For this reason, a collective oscillation whose frequency lies below this interval is not in resonance with any gas particle and hence cannot be damped by mixing.

The question of whether small perturbations of a steady state of the gravitational Vlasov-Poisson system are damped or give rise to persistent oscillations goes back to Lynden-Bell~\cite{dL62}. Radial oscillations which persist over many dynamical times have been reported in numerical experiments with spherical shells and in $N$-body simulations since the work of H\'enon, see Refs.~\cite{mH73,pLoG88} and references therein. The theoretical analysis started with the one-dimensional system of plane-parallel mass sheets, which serves as a model for the vertical structure of a galactic disk and whose steady states of finite extent likewise possess a band of orbital frequencies. Mathur~\cite{sM90} reduced the existence of an oscillating mode whose frequency lies in a gap of the continuous spectrum, which consists of this band and its harmonics, to the condition that a symmetric integral operator has the eigenvalue one. His analysis covers the one-dimensional system and the radial perturbations of spherical systems, in both cases in the presence of a fixed external field, and it shows that such modes exist for certain steady states. For the one-dimensional system, Weinberg~\cite{mW91} located these modes numerically for families of truncated isothermal and polytropic slabs, with and without an external halo, and reproduced them in $N$-body simulations. By means of a fluid model and $N$-body simulations, Louis~\cite{pL92} determined the members of a sequence of polytropes and of King models at which the modes enter the band and turn into weakly damped density waves; the damped oscillations of truncated and untruncated isothermal slabs have been computed in Ref.~\cite{lWgB15}. A closely related problem arises in the physics of particle accelerators. There, the longitudinal oscillations of a bunch in the potential well of the radio-frequency system lose their Landau damping when a coherent mode emerges from the band of the synchrotron frequencies, and the corresponding intensity thresholds are obtained from the Vlasov equation linearized in action-angle variables, see Ref.~\cite{iKtAeS21} and references therein.

For spherically symmetric steady states the problem has received considerable attention in recent years. Had\v{z}i\'c, Rein and Straub~\cite{mHgRcS22} (see also Refs.~\cite{Kunze-Book,mK22}) showed that the essential spectrum of the operator governing the linearized dynamics about a large class of spherically symmetric steady states is determined by the radial periods of the gas particles. Using a version of the Birman-Schwinger principle which was first applied to this problem by Mathur~\cite{sM90}, they derived a criterion for the existence of an eigenvalue in the gap below this spectrum, which corresponds to an oscillating mode that is not damped. This criterion is formulated in terms of a family of bounded operators acting on functions of the radius. Related criteria for the existence of oscillating modes and estimates for their number, which also allow for an attractive external potential, have been given by Moreno, Rioseco and Van Den Bosch~\cite{mMpRhV26} for steady states whose DF depends only on the energy. For the reduced system in which all the gas particles have the same magnitude of the angular momentum and move around a central point mass, Had\v{z}i\'c, Rein, Schrecker and Straub~\cite{mHgRmScS25} proved the following dichotomy for steady states of sufficiently small mass whose DF is proportional to $(E_0 - E)^k$ for energies $E$ below the cutoff energy $E_0$: the solutions of the linearized system are damped if $k > 1$, whereas an oscillating mode exists if $1/2 < k\leq 1$. Quantitative decay rates for the linearized system about steady states with a central point mass and a distribution of angular momenta have been obtained in Ref.~\cite{mHmS25}, and for the nonlinear system with an external Kepler potential, phase mixing of spherically symmetric data over long time scales has been established in Ref.~\cite{sCjL26}. On the numerical side, Ramming and Rein~\cite{tRgR18} and Straub~\cite{cS24} have performed particle-in-cell simulations of the spherically symmetric Vlasov-Poisson system which show that, depending on the steady state, its perturbations either oscillate without damping or decay. Finally, we mention that weakly damped non-radial modes of spherical stellar systems have been computed in Refs.~\cite{mW94,jFsP22} by means of the matrix method, that the frequencies of the lowest radial modes of a family of spherical models have been obtained with a related method in Ref.~\cite{pV04}, and that basis functions adapted to the radial perturbations of systems of infinite extent have recently been constructed in Ref.~\cite{ePiS26}.

In these notes, we study numerically the transition between mixing and persistent oscillations for a spherically symmetric collisionless gas which is trapped in an external central potential and whose self-gravity is taken into account. We consider two situations. The first one is the reduced model in which all the gas particles have the same magnitude $L_0 > 0$ of the angular momentum, which has a two-dimensional phase space and is the setting considered in Refs.~\cite{pRoS20,mHgRmScS25}. This model goes back to the numerical experiments of H\'enon with spherical shells (see Ref.~\cite{mH73} and references therein) and has been used in the $N$-body simulations of Ref.~\cite{pKbM02}; moreover, the example for which Mathur~\cite{sM90} established the existence of a radial mode is a thin shell of nearly circular orbits with a narrow distribution of angular momenta in an external field. In the second situation the angular momenta of the gas particles are distributed over an interval. For the external potential we choose the isochrone potential, for which the action-angle variables are known explicitly, and the potential generated by a point mass. Our analysis is based on three independent tools: (i) a particle-in-cell (PIC) code which solves the nonlinear Vlasov-Poisson system, (ii) a numerical solution of the equations linearized about a self-consistent steady state, formulated in the action-angle variables of the latter, and (iii) the numerical computation of the largest eigenvalue $\lambda(\omega)$ of the operator of Mathur, which measures the gain of the self-consistent loop in which a perturbation of the potential with frequency $\omega$ induces a perturbation of the density which, in turn, generates a potential. An oscillating mode with frequency below the frequency band of the gas particles exists if and only if the value $\lambda_{edge}$ of this gain at the lower edge of the band is larger than one. This is the criterion of Refs.~\cite{sM90,mHgRcS22,mK22}; the quantity $\lambda_{edge}$ corresponds to the number which is denoted by $M$ in Ref.~\cite{mHgRcS22} and by $\mu_*$ in Ref.~\cite{mK22}, and its determination is listed among the open questions in the latter reference.

The main results of these notes can be summarized as follows. First, as long as the mass of the gas is sufficiently small compared to the mass which generates the external potential, a small perturbation of a self-consistent steady state decays, and the response obtained from the PIC simulations agrees with the solution of the linearized equations to within one percent. Close to the threshold mentioned below the decay is exponential, as in Landau damping, and at late times it is algebraic. Second, for the reduced model there is a threshold mass above which a discrete mode emerges below the frequency band, such that the perturbation oscillates without being damped. We show that this threshold is characterized by the condition $\lambda_{edge} = 1$, and that the frequency of the mode predicted by the equation $\lambda(\omega) = 1$ agrees with the one measured in the simulations. The corresponding phenomenon for the one-dimensional slab was described in Refs.~\cite{mW91,pL92}. Here we determine the threshold masses for the radial problem with a central potential, in particular for the polytropes of Ref.~\cite{mHgRmScS25}, whose results hold below a mass which is not estimated in that reference, and we obtain the laws which govern the gain and the frequency of the mode close to the edge of the band as functions of the exponent $k$. Third, for steady states of small mass with $k\leq 1$, the oscillating mode whose existence is guaranteed by the results in Ref.~\cite{mHgRmScS25} lies so close to the edge of the band that it cannot be distinguished from the decaying part of the response over the time span of our simulations; instead, for all the values of $k$ we have considered the response decays algebraically, and over the time span of our simulations it follows the law $t^{-(k+1)}$ of the mixing in a fixed potential. We also find that at later times, or for larger masses, the self-gravity of the perturbation reduces the exponent of this decay from $k+1$ to $k$, which is the exponent of the decay estimates of Ref.~\cite{mHgRmScS24} for the mixing, in a fixed potential, of perturbations that vanish at the edge like the derivative of the DF of the steady state. Fourth, the oscillating modes which lie close to the edge of the band are very sensitive to the amplitude of the perturbation: above a critical amplitude, which we estimate by means of a pendulum model and which vanishes at the threshold, the gas particles near the edge of the steady state are trapped by the mode, and the amplitude of the oscillation departs from the one predicted by the linearized equations. That the amplitude which such a mode can sustain is limited by its closeness to the band was observed by Weinberg~\cite{mW91} in simulations of the one-dimensional slab; the pendulum model provides an explicit estimate for the amplitude at which this limitation sets in. Fifth, a dispersion in the angular momentum removes the divergence of $\lambda(\omega)$ at the edge of the band which is responsible for the oscillating modes with $k\leq 1$ in the reduced model, such that these modes disappear when the mass of the gas is small enough, and it suppresses the coherent response which is present when all the gas particles have the same angular momentum.\EJ{The last statement relies on the operator with dispersion and on the free solution; the PIC runs with dispersion of section~\ref{Sec:Dispersion} are still to be done.}

Particle methods have a long tradition in the numerical study of spherically symmetric collisionless systems, both in the Newtonian case~\cite{tRgR18,cS24} and in the relativistic one, where they have been used to analyze the collapse of star clusters and the critical behavior at the threshold of black hole formation~\cite{sStT85,gRaRjS98,Olabarrieta:2000bn,Olabarrieta:2001wy,hAgR06}; see Refs.~\cite{BirdsallLangdon-Book,HockneyEastwood-Book,lapenta2015kinetic} for introductions to the PIC method, Ref.~\cite{Shalaby:2017yok} for a higher-order implementation in the context of plasma physics, and Ref.~\cite{gRtR03} for a convergence proof of a PIC scheme for the spherically symmetric Vlasov-Einstein system. Measuring the response of the gas to a small perturbation with such a method is delicate, because the signal has to be separated from the noise which is due to the finite number of computational particles. In these notes we address this problem by placing the particles on a regular lattice in the action-angle variables of the self-consistent steady state, by encoding the perturbation in the weights of the particles instead of their positions (as in Refs.~\cite{fLfCjB93,jBaOyY11}), and by subtracting a reference run without perturbation. A loading of the particles which reduces the initial fluctuations (a quiet start) was already used for systems of spherical shells in Ref.~\cite{mH73}. We also show that the observables need to be computed with the action-angle variables belonging to the total potential, since those of the external potential alone leave a time-independent residue which could be mistaken for an undamped mode.

The remainder of these notes is organized as follows. In section~\ref{Sec:VPSphSym} we formulate the Vlasov-Poisson system with an external potential in spherical symmetry, for the reduced model with fixed angular momentum as well as for a distribution of angular momenta, and we introduce the action-angle variables and the observables used throughout these notes. In section~\ref{Sec:Linear} we review the mixing phenomenon in the absence of self-gravity, derive the linearized equations about a steady state and introduce the gain $\lambda(\omega)$ of the self-consistent loop. Our numerical methods are described in section~\ref{Sec:Numerics}, and their validation is presented in section~\ref{Sec:Validation}. The results for the reduced model are discussed in section~\ref{Sec:FixedL}, and those for a distribution of angular momenta in section~\ref{Sec:Dispersion}. Conclusions and an outlook to future work are given in section~\ref{Sec:Conclusions}. Technical details regarding the numerical computation of the action-angle variables are included in Appendix~\ref{App:ActionAngle}.



%%%%%%%%%%%%%%%%%%%%%%%%%%%%%%%%%%%%%%%%%%%%%
\section{The Vlasov-Poisson system in spherical symmetry}
\label{Sec:VPSphSym}
%%%%%%%%%%%%%%%%%%%%%%%%%%%%%%%%%%%%%%%%%%%%%

\EJ{To be written: model, reduction to spherical symmetry with fixed and with dispersed angular momentum, units, action-angle variables, observables.}


%%%%%%%%%%%%%%%%%%%%%%%%%%%%%%%%%%%%%%%%%%%%%
\section{Mixing and linear response}
\label{Sec:Linear}
%%%%%%%%%%%%%%%%%%%%%%%%%%%%%%%%%%%%%%%%%%%%%

\EJ{To be written: mixing without self-gravity, linearized equations, gain of the self-consistent loop, the case with a distribution of angular momenta.}


%%%%%%%%%%%%%%%%%%%%%%%%%%%%%%%%%%%%%%%%%%%%%
\section{Numerical methods}
\label{Sec:Numerics}
%%%%%%%%%%%%%%%%%%%%%%%%%%%%%%%%%%%%%%%%%%%%%

\EJ{To be written: particle-in-cell method, equilibria and initial data, linearized solver, evaluation of the gain.}


%%%%%%%%%%%%%%%%%%%%%%%%%%%%%%%%%%%%%%%%%%%%%
\section{Code validation}
\label{Sec:Validation}
%%%%%%%%%%%%%%%%%%%%%%%%%%%%%%%%%%%%%%%%%%%%%

\EJ{To be written.}


%%%%%%%%%%%%%%%%%%%%%%%%%%%%%%%%%%%%%%%%%%%%%
\section{Results for fixed angular momentum}
\label{Sec:FixedL}
%%%%%%%%%%%%%%%%%%%%%%%%%%%%%%%%%%%%%%%%%%%%%

\EJ{To be written.}


%%%%%%%%%%%%%%%%%%%%%%%%%%%%%%%%%%%%%%%%%%%%%
\section{Results for a distribution of angular momenta}
\label{Sec:Dispersion}
%%%%%%%%%%%%%%%%%%%%%%%%%%%%%%%%%%%%%%%%%%%%%

\EJ{To be written.}


%%%%%%%%%%%%%%%%%%%%%%%%%%%%%%%%%%%%%%%%%%%%%
\section{Conclusions}
\label{Sec:Conclusions}
%%%%%%%%%%%%%%%%%%%%%%%%%%%%%%%%%%%%%%%%%%%%%

\EJ{To be written.}


%%%%%%%%%%%%%%%%%%%%%%%%%%%
%%%   ACKNOWLEDGMENTS
%%%%%%%%%%%%%%%%%%%%%%%%%%%

\acknowledgments

\EJ{To be written.}


\appendix

%%%%%%%%%%%%%%%%%%%%%%%%%%%%%%%%%%%%%%%%%%%%%
\section{Numerical computation of the action-angle variables}
\label{App:ActionAngle}
%%%%%%%%%%%%%%%%%%%%%%%%%%%%%%%%%%%%%%%%%%%%%

\EJ{To be written.}


\bibliographystyle{unsrt}
\bibliography{refs}

\end{document}
